\section{Síntesis y ampliación}

La semana 9 cambia la perspectiva de «escribir código» a «operar código». DevOps y SRE son la parte que más tiempo consume en la ingeniería de software, y AI SRE está reduciendo la carga operativa mediante grafos de conocimiento, agentes que abarcan toda la pila y análisis en tiempo real.

\subsection{Puntos clave}

\begin{itemize}
    \item Escribir código representa solo el 30\% del tiempo de ingeniería; la operación es el 70\% restante, más difícil
    \item Cuatro señales de oro: latencia, errores, tráfico, saturación
    \item La respuesta a incidentes sigue un playbook estandarizado; la métrica central es el MTTR
    \item Cuatro características de AI SRE: grafo de conocimiento dinámico, agentes que abarcan toda la pila, narración en tiempo real, interpretabilidad
    \item El mayor valor: romper los silos de conocimiento y reducir el MTTR
    \item Actualmente se empieza por el análisis de causa raíz; la reparación automática aún está en una etapa temprana
\end{itemize}

\subsection{Lista de implementación operativa}

\begin{enumerate}
    \item Primero unificar los criterios de monitoreo y la calidad de los registros, y solo después introducir AI SRE; no dejar que la IA sustituya una base de datos desordenada
    \item Tomar el MTTR, la tasa de falsos positivos y la precisión de escalamiento como métricas centrales, en lugar de fijarse solo en si la respuesta del modelo «parece» adecuada
    \item Automatizar dentro de un alcance controlado las operaciones de bajo riesgo, y dejar las reparaciones de alto riesgo a la aprobación humana
\end{enumerate}

\subsection{Tabla de reparto de roles en guardia}

\begin{table}[H]
\centering
\begin{tabular}{p{0.18\textwidth} p{0.26\textwidth} p{0.3\textwidth}}
\toprule
Rol & Responsabilidad en el sistema AI SRE & Responsabilidad humana que no puede delegarse \\
\midrule
Ingeniero de guardia & Revisar el resumen de la IA, decidir la acción de mitigación, coordinar el escalamiento & Responsable de la confirmación final de las operaciones importantes \\
Sistema AI SRE & Agregar evidencia, proponer hipótesis de causa raíz, sugerir acciones de bajo riesgo & No puede ejecutar reparaciones de alto riesgo sin autorización \\
Equipo de plataforma/infraestructura & Mantener el monitoreo, los registros, los runbooks y los límites de permisos & Garantizar que la IA disponga de datos operativos de alta calidad \\
\bottomrule
\end{tabular}
\caption{Límites entre roles humanos y de máquina en el escenario AI SRE}
\end{table}

\subsection{Por qué el ciclo de postmortem es aún más importante}

AI SRE no hace que el postmortem de incidentes sea menos importante; al contrario, lo vuelve más importante. Esto es porque cada atención de un incidente expone al mismo tiempo dos tipos de problemas:

\begin{itemize}
    \item Problemas del sistema mismo: vacíos de monitoreo, dependencias poco claras, runbooks obsoletos, límites de permisos ambiguos
    \item Problemas del sistema de IA: qué señales pasó por alto, qué ruido amplificó, en qué paso dio una recomendación poco confiable
\end{itemize}

Por lo tanto, los equipos maduros necesitan reescribir los resultados del postmortem en el monitoreo, las reglas de alerta, los prompts de la IA, los runbooks y los límites de aprobación, para que AI SRE no sea un asistente de incidentes de un solo uso, sino un sistema operativo que evoluciona de manera sostenida.

\subsection{Lecturas adicionales}

\begin{itemize}
    \item Resolve AI: \url{https://resolve.ai}
    \item Google SRE Book: \url{https://sre.google/sre-book/table-of-contents/}
    \item DataDog Bits AI: \url{https://www.datadoghq.com/product/platform/bits-ai/}
    \item ``The Four Golden Signals'' (Google SRE Book, Chapter 6)
    \item PagerDuty Incident Response Guide: \url{https://response.pagerduty.com}
\end{itemize}

%% --- Fin del contenido --- %%

\end{document}
